% Options for packages loaded elsewhere
\PassOptionsToPackage{unicode}{hyperref}
\PassOptionsToPackage{hyphens}{url}
\documentclass[
  ngerman,
]{article}
\usepackage{xcolor}
\usepackage{amsmath,amssymb}
\setcounter{secnumdepth}{-\maxdimen} % remove section numbering
\usepackage{iftex}
\ifPDFTeX
  \usepackage[T1]{fontenc}
  \usepackage[utf8]{inputenc}
  \usepackage{textcomp} % provide euro and other symbols
\else % if luatex or xetex
  \usepackage{unicode-math} % this also loads fontspec
  \defaultfontfeatures{Scale=MatchLowercase}
  \defaultfontfeatures[\rmfamily]{Ligatures=TeX,Scale=1}
\fi
\usepackage{lmodern}
\ifPDFTeX\else
  % xetex/luatex font selection
\fi
% Use upquote if available, for straight quotes in verbatim environments
\IfFileExists{upquote.sty}{\usepackage{upquote}}{}
\IfFileExists{microtype.sty}{% use microtype if available
  \usepackage[]{microtype}
  \UseMicrotypeSet[protrusion]{basicmath} % disable protrusion for tt fonts
}{}
\makeatletter
\@ifundefined{KOMAClassName}{% if non-KOMA class
  \IfFileExists{parskip.sty}{%
    \usepackage{parskip}
  }{% else
    \setlength{\parindent}{0pt}
    \setlength{\parskip}{6pt plus 2pt minus 1pt}}
}{% if KOMA class
  \KOMAoptions{parskip=half}}
\makeatother
\ifLuaTeX
\usepackage[bidi=basic,shorthands=off,]{babel}
\else
\usepackage[bidi=default,shorthands=off,]{babel}
\fi
\ifLuaTeX
  \usepackage{selnolig} % disable illegal ligatures
\fi
\setlength{\emergencystretch}{3em} % prevent overfull lines
\providecommand{\tightlist}{%
  \setlength{\itemsep}{0pt}\setlength{\parskip}{0pt}}
\usepackage{bookmark}
\IfFileExists{xurl.sty}{\usepackage{xurl}}{} % add URL line breaks if available
\urlstyle{same}
\hypersetup{
  pdftitle={Drei Filme Zwei B├╝cher},
  pdfauthor={Paul Koop},
  pdflang={de},
  hidelinks,
  pdfcreator={LaTeX via pandoc}}

\title{Drei Filme Zwei B├╝cher}
\usepackage{etoolbox}
\makeatletter
\providecommand{\subtitle}[1]{% add subtitle to \maketitle
  \apptocmd{\@title}{\par {\large #1 \par}}{}{}
}
\makeatother
\subtitle{Zwei späte Romane Don DeLillos, zwei Filme von Jaco Van
Dormael und ein Film, den es so nie gegeben hat}
\author{Paul Koop}
\date{}

\begin{document}
\maketitle

{
\setcounter{tocdepth}{3}
\tableofcontents
}
\section{Vorbemerkung}\label{vorbemerkung}

Dieser Essay ist kein Kommentar zu den Erzählungen des Pompeji-Projekts.
Er ist ein eigenständiger Zugang zu denselben Fragen, die die Romane und
Erzählungen narrativ entfalten. Er setzt sie nicht voraus. Er erklärt
sie nicht. Er ist eine andere Stimme im selben Gespräch, eine, die
beschreibt, was die Erzählungen vorführen.

Die folgenden Analysen kreisen um zwei Werke Don DeLillos, zwei Filme
Jaco Van Dormaels und einen Film, den es so nie gegeben hat. Sie tun das
nicht, um diese Werke zu bewerten oder zu vereinnahmen. Sie tun es, um
eine Denkbewegung sichtbar zu machen, die sich in ihnen artikuliert,
eine Bewegung, die auch das Pompeji-Projekt trägt. Wer die Theorie
sucht, wird sie hier finden. Wer die Erzählung sucht, wird anderswo
fündig. Beide Wege führen zu demselben Rand, aber keiner setzt den
anderen voraus.

\section{Vorspiel}\label{vorspiel}

Es geht in den folgenden Texten nicht um Bücher und Filme im üblichen
Sinn. Nicht um Bewertungen, nicht um Empfehlungen, nicht um
abgeschlossene Urteile. Was hier nebeneinandergestellt wird, zwei späte
Romane Don DeLillos, zwei Filme von Jaco Van Dormael und ein Film, den
es so nie gegeben hat -, bildet ein lose gefügtes Ensemble. Kein Kanon,
kein System. Eher ein Resonanzraum. Literatur und Film werden dabei
nicht getrennt betrachtet, sondern als unterschiedliche Ausdrucksformen
derselben Frage gelesen: Was bleibt, wenn die großen Vermittlungen
verstummen? Alle Werke kreisen um ähnliche Motive, ohne je dasselbe zu
sagen. Um Stille, nicht als Idylle, sondern als Zustand nach dem Lärm.
Um Zeit, nicht als Uhrzeit, sondern als Möglichkeit, Verdichtung,
Rücklauf. Um Ordnung, technische, religiöse, moralische -, die Sinn
verspricht und Macht ausübt. Und um das, was sich dieser Ordnung
entzieht: Möglichkeit, Beziehung, Offenheit. In DeLillos Roman Der
Omega-Punkt führt die Stille nicht in Leere, sondern in eine
eigentümliche Zeitlosigkeit, in der Erkenntnis total und nicht mehr
sagbar ist, vielleicht aber erfahrbar. In Die Stille fällt dann die
Technik aus, und mit ihr verschwinden Sprache, Deutung, Gewissheit.
Zurück bleibt eine Gegenwart, die nichts erklärt und nichts erlöst. Zwei
Formen der Stille stehen einander gegenüber: die eine als metaphysische
Verdichtung, die andere als reiner Entzug. Van Dormaels Filme öffnen
einen anderen Zugang. In Mr. Nobody zerfällt das Leben in Möglichkeiten.
Jede Entscheidung erzeugt eine Wirklichkeit, jede Geschichte ist real,
weil sie möglich ist. Zeit ist kein Strom, sondern ein Feld. In Das
brandneue Testament wird ein Gottesbild entmachtet, das aus Angst,
Schuld und Kontrolle besteht. Nicht durch Widerlegung, sondern dadurch,
dass es seine Macht verliert. Sinn entsteht hier nicht aus Ordnung,
sondern aus Beziehung. Der fiktive Film Gott wohnt im dritten Stock
treibt diese Bewegung weiter. Gott erscheint nicht mehr als
transzendente Instanz, sondern als internalisiertes Regime aus Schuld
und Gewissheit. Erlösung geschieht nicht durch neue Lehre, sondern durch
das Verlassen dieses Systems. Die Figuren, die tragen, sind keine
Heiligen und keine Lehrer, sondern Menschen ohne metaphysischen
Anspruch, tragfähig, weil sie nichts rechtfertigen müssen. Was all diese
Texte verbindet, ist keine Botschaft, sondern eine Bewegung. Sie führen
nicht zu einer neuen Synthese, keinem neuen Bund, keiner letzten
Wahrheit. Sie zeigen vielmehr, wie Macht verschwindet: die Macht der
Technik, des allwissenden Erzählers, des strafenden Gottes, der großen
Sinnnarrative. Nicht durch Zerstörung, sondern durch Irrelevanz. Zurück
bleibt kein Nihilismus. Sondern eine fragile Form von Humanität. Still,
vorläufig, ohne Garantie. Eine Existenz, die Leid, Zeit und Endlichkeit
nicht aufhebt, aber auch nicht mehr rechtfertigen muss, vielleicht der
Omegapunkt. Mit diesem Hintergrund lesen sich die folgenden Texte nicht
als Einzelrezensionen, sondern als Variationen eines Gedankens. Als
Annäherungen an eine Welt, in der Sinn nicht verschwindet, aber seine
Aussprechbarkeit verliert, vielleicht im Omegapunkt. In den Werken
DeLillos und Van Dormaels artikuliert sich ein unausgesprochenes
katholisches Lebensgefühl, das Unvollkommenheit, Leid und Zeit nicht
rechtfertigt, sondern annimmt. Dieses Lebensgefühl entzieht einem
augustinisch-calvinischen, dualistischen Gottesbild die Grundlage, nicht
durch Widerlegung, sondern durch Bedeutungsverlust. Heilung erscheint
nicht als Erlösung von der Welt, sondern als zeitlose Vereinigung
innerhalb des Gelebten.

\section{Zwei Stimmen, ein Grenzwert}\label{zwei-stimmen-ein-grenzwert}

Es wäre ein Missverständnis, diese Texte als Variationen eines Gedankens
zu lesen. Sie sind Variationen zweier Gedanken, die sich nicht
vereinbaren lassen. Der eine ist der Gedanke DeLillos: dass die Stille
nicht in die Tiefe führt, sondern in die Ausdünnung. Dass der
Omega-Punkt kein Ziel ist, sondern ein Grundzustand, der Wärmetod, die
Entropie, das ungemessene Material. Der andere ist der Gedanke Van
Dormaels: dass die Stille nicht das Ende ist, sondern die Fülle. Dass
der Omega-Punkt kein Grundzustand ist, sondern ein Feld, die Komplexität
der Möglichkeiten, die Differenz, die nicht nivelliert wird.

Diese beiden Gedanken lassen sich nicht zu einer Synthese bringen. Wer
sie vereinbaren will, zerstört sie. Wer einen von beiden für den ganzen
hält, macht ihn zur Ideologie. Teilhard de Chardin hat genau das getan:
Er hat den Gedanken der Richtung, der Teleologie, für den Omega-Punkt
selbst gehalten. Er war Augustiner. Er war Thomist. Er konnte nicht
anders, als die Zeit auf ein Ende hin zu denken. Das war seine Größe,
und sein Irrtum.

Der Omega-Punkt ist nicht das Ziel. Er ist nicht das Ende. Er ist der
Rand, an dem beide Gedanken gleichzeitig gelten, ohne sich zu versöhnen.
Die folgenden Analysen versuchen, diesen Rand sichtbar zu machen. Sie
führen nicht zu einer Antwort. Sie führen zu der Frage, die beide
Register offen lassen.

\section{Der Omega-Punkt}\label{der-omega-punkt}

DeLillo, D.: Der Omega-Punkt. ISBN-10: 3462041924, ISBN-13:
978-3462041927. Übersetzung von Heibert, F. Original: Point Omega, 2010.

Autor: Der Autor wuchs in der Zwischenkriegszeit, wie Jaco Van Dormael,
in einer katholischen Familie mit italienischen Wurzeln in der Bronx
auf. Spät findet er zur Literatur. Er arbeitet als Texter mit
Bachelorabschluss in der Werbebranche und beginnt zu schreiben. Seine
Werke werden weltweit bekannt und kreisen um die postmoderne Reflexion
der Welt, des Lebens und der krisenhaften Entwicklung in den USA.

Inhalt: Zwei Männer betrachten im ersten Teil im Museum of Modern Arts,
New York, die Videoinstallation 24 Hour Psycho von Douglas Gordon, in
der Alfred Hitchcocks Spielfilm Psycho auf eine Laufzeit von 24 Stunden
gedehnt abgespielt wird. Im nächsten Teil des Romans folgt ein junger
Filmemacher einem aus dem Dienst ausgeschiedenen Kriegsberater der
Bush-Regierung in die Wüste, um über ihn eine Dokumentation zu fertigen.
Der alte Mann ist in die Wüste gegangen, um der Zeitlosigkeit zu folgen.
Hinter dem Leben und nach dem Bewusstsein sieht der alte Mann den
Omegapunkt Teilhard de Chardins. Noch bevor der Filmemacher mit seiner
Dokumentation beginnen kann, kommt die Tochter des alten Mannes. Die
Mutter hat sie geschickt, damit sie in der Wüste über die Absicht eines
Verehrers nachdenkt, mit ihr eine Lebenspartnerschaft einzugehen. Kaum
hat man sich kennengelernt, verschwindet die junge Frau und bleibt
verschwunden.

Wertung: Das Thema des Romans ist die Stille, nicht die
Geräuschlosigkeit, sondern die Stille. Die Stille ist Zeitlosigkeit. Es
ist die Zeitlosigkeit der Steine, in der die Zwischenstufe Leben sich
verflüchtigt. Und diese Zeitlosigkeit der Steine ist das Tor zu einem
Ozean des Wissens. Dieses Wissen erschließt sich, wenn man das Leben
durchlitten hat. Dann weiß man, was Zeitlosigkeit, Liebe und Leid ist
und kann es doch nicht mehr sagen. Die Zeitlosigkeit steckt in dem zum
Standbild verlangsamten Moment der Filmsequenz. Die Stille der Steine
steckt im Rückzug in die Wüste zu jesuitischen Exerzitien und der
Verlust, die Liebe und das Leid stecken im Verschwinden der Tochter.

\section{Die Stille}\label{die-stille}

DeLillo, D.: Die Stille. ISBN: 978-3-462-30295-0. Übersetzung von
Heibert, F. Original: The Silence, 2020.

Autor: Don DeLillo, geboren 1936 in der Bronx, der wie Jaco Van Dormael
in einem katholischen Umfeld aufgewachsen ist, stammt aus einer
katholischen Familie italienischer Herkunft. Er studiert Kommunikation
und arbeitet zunächst in der Werbung, bevor er sich spät, aber
entschieden der Literatur zuwendet. Seit den 1970er-Jahren entwickelt er
sich zu einem der bedeutendsten Chronisten der Spätmoderne. Seine Romane
kreisen um das Bewusstsein in technologisch beschleunigten Zeiten, die
Unsichtbarkeit von Macht, die Auflösung individueller Identität und die
Frage, was Sprache noch retten kann.

Inhalt: Am Tag des Super Bowl, mitten in einem New Yorker Wohnzimmer,
fallen sämtliche elektronischen Systeme aus. Bildschirme erlöschen. Kein
Fernsehen, kein Telefon, kein Internet. Eine kleine Gruppe von Menschen,
ein Ehepaar, ein Freund, ein ehemaliger Physikprofessor und dessen Frau,
sitzt zusammen, während draußen eine Welt im technischen Blackout
versinkt. Der Roman beginnt mit einer prekären Flugzeuglandung und endet
in einer fast rituellen Zusammenkunft bei Kerzenschein. Die Gespräche
taumeln zwischen Alltagsbanalitäten, Halbsätzen über Einstein,
Sonnenstürmen und digitalen Kontrollregimen. Die Worte werden weniger,
die Pausen länger. Jeder Rückgriff auf eine Erklärung, Sonnenflecken,
Systemversagen, führt tiefer in eine Leere hinein. Die Handlung bleibt
fragmentarisch, fast statisch. Aber gerade darin liegt das Programm des
Romans: In der Stille, in der plötzlichen Entelechie des Schweigens,
enthüllt sich die Fragilität des modernen Lebens. Was geschieht mit uns,
wenn der Informationsfluss endet?

Wertung: Die Stille ist kein Katastrophenroman. Es ist ein meditativer
Text über die Abwesenheit von Datenströmen, über die Entblößung des
Menschlichen im Moment der medialen Entwöhnung. Der Titel meint nicht
bloß das Verstummen der Technik, sondern das Verstummen der Deutungen.
DeLillo legt eine stille Miniatur vor, eine Art literarische Meditation
über das Ende der Vermittlung. In der Stille offenbart sich nicht
Erlösung, sondern das Verstummen der Gewissheiten. Und doch birgt diese
Stille eine Ahnung von Transzendenz, ein leiser Nachhall metaphysischer
Fragen, der durch Einsteins Satz „Ich weiß nicht, mit welchen Waffen der
dritte Weltkrieg ausgetragen wird..." hindurchtönt. Der Text ist
minimalistisch, fast theaterhaft. Die Charaktere sprechen wie aus dem
Off, verlangsamt, wie in einem inneren Echo. Jeder scheint mit sich
selbst zu sprechen, auch wenn andere im Raum sind. Und was bleibt, ist
die Erinnerung an eine Welt, die sich für unersetzbar hielt, bis sie
verschwand.

\#\#\# Vergleich mit Der Omegapunkt

Während Der Omegapunkt von metaphysischer Dichte getragen ist, ein
Roman, der mit Teilhard de Chardins Visionen flirtet, die Wüste als Raum
der Erkenntnis erforscht und das Verstummen als Eintritt in die Ewigkeit
deutet -, bleibt Die Stille auf eine radikal gegenwärtige Leere
fokussiert. Im Omegapunkt führt das Verstummen in die Tiefe, in eine
mystisch aufgeladene Sphäre jenseits der Zeit. Die Wüste steht dort für
das Göttliche im Rückzug, für die langsame Auslöschung des Egos in einer
kontemplativen Ekstase. Die Stille ist hier der Auftakt zu Erkenntnis.
In Die Stille hingegen ist das Schweigen nicht Erlösung, sondern Entzug.
Es bleibt keine göttliche Sphäre, kein metaphysisches Ziel, sondern ein
Wohnzimmer mit Menschen, die von Einsteins Sprache träumen, aber nicht
mehr sprechen können. Der Omegapunkt führt in eine jenseitige Ordnung;
Die Stille konfrontiert uns mit dem Verschwinden der Ordnung überhaupt.

\section{Zwei DeLillo-Stimmen}\label{zwei-delillo-stimmen}

Es lohnt sich, die beiden DeLillo-Bücher nicht nur gegeneinander,
sondern als zwei Stimmen innerhalb desselben Denkens zu lesen. Der
Omega-Punkt tendiert zur metaphysischen Tiefe: Die Stille ist
Zeitlosigkeit, die Wüste ist ein Raum der Erkenntnis, das Verschwinden
der Tochter ist ein Ereignis, das über das Kriminalistische hinausweist.
Die Stille tendiert zur radikalen Leere: Das Verstummen der Technik
führt nicht in die Tiefe, sondern in das Wohnzimmer, in dem Menschen
nicht mehr sprechen können. Der eine Text lässt die Transzendenz
durchscheinen. Der andere lässt sie verschwinden.

Beide Stimmen gehören zu DeLillo. Aber sie gehören nicht zusammen.
Zwischen ihnen liegt kein Übergang, keine Entwicklung, keine Versöhnung.
Sie sind zwei Zustände desselben Denkens, und das Denken selbst weiß
nicht, welcher Zustand der wahre ist. Genau darin liegt die
DeLillo-Perspektive auf den Omega-Punkt: Sie ist nicht einseitig. Sie
ist innerlich gespalten. Wer sie auf einen Ton festlegt, sei es die
Tiefe oder die Leere -, verliert genau das, was sie ausmacht.

Fazit: Die Stille ist ein Spätwerk, das wie ein literarischer Abspann
wirkt. Es ist DeLillos kürzestes und wohl auch kargstes Buch. Sprachlich
konzentriert, thematisch entkernt, formbewusst leer. Es ist nicht die
Antwort auf eine Krise, sondern die Frage, ob wir überhaupt noch auf
Krisen antworten können, ohne Netz, ohne Bildschirm, ohne Worte.

Wer im Omegapunkt noch Spuren einer Offenbarung entdeckt hat, findet
hier nur noch das Echo einer untergegangenen Welt. Und vielleicht genau
darin liegt die stille Radikalität dieses Textes.

\section{Mr. Nobody}\label{mr.-nobody}

Mr. Nobody von Jaco Van Dormael, 2009.

Autor: Jaco Van Dormael ist ein belgischer Regisseur und Drehbuchautor,
der wie Don DeLillo in einem katholisch geprägten Umfeld aufwuchs und
der für seine innovativen und visuell beeindruckenden Filme bekannt ist.
Geboren am 9. Februar 1957 in Ixelles, Belgien, lernte und arbeitete er
in den Bereichen Film und Theater. Van Dormael wurde international
bekannt durch seine Filme. Er zeichnet sich durch seinen handwerklich
beeindruckenden, einzigartigen Stil aus, der eine Mischung aus
Surrealismus, Romantik und philosophischen Themen verkörpert.

Handlung: Der Film Mr. Nobody beginnt mit einer symbolischen Szene, in
der eine Taube in einer Skinnerbox trainiert wird. Diese Szene
symbolisiert das Konzept der Konditionierung und wie wir oft glauben,
dass unser zufälliges Verhalten ursächlich für Ereignisse ist, was oft
eine Illusion darstellt.

In den folgenden Einstellungen taucht eine mit sich identische Person
unkommentiert in voneinander unabhängigen alternativen Geschichten auf.

Der Zuschauer erfährt erst dann, dass diese Person Mr. Nobody der letzte
sterbliche Mensch ist und sein bevorstehender Tod live in den Medien
übertragen wird, begleitet von einem Psychiater und später einem
Journalisten.

Der Psychiater führt Mr. Nobody durch Hypnose in die Zeit vor seiner
Geburt zurück. Hier wird enthüllt, dass vor der Geburt alles Wissen über
mögliche Entwicklungen und Lebenswege vorhanden ist, aber die
sogenannten „Engel des Vergessens" nehmen dieses Wissen bei der Geburt
weg. Bei Mr. Nobody haben die Engel vergessen, es zu nehmen, was ihm
ermöglicht, sich an alle möglichen Leben zu erinnern und auch die
zukünftigen Leben zu kennen.

Die Geschichte konzentriert sich auf einen jungen Jungen, der bei der
Trennung seiner Eltern vor der Wahl steht, zu seinem Vater oder seiner
Mutter zu gehen. Er erinnert sich an drei Geschichten mit
unterschiedlichen Ehepartnern und Kindern, sowie an eine
Sci-Fi-Geschichte über die Marsbesiedelung, die er als Jugendlicher
verfasst hat. Diese Geschichte ist ebenso real wie die anderen, da sie
möglich ist.

Im Laufe seiner möglichen Leben ist Mr. Nobody auch als Fernsehmoderator
tätig und behandelt in seiner Dokumentationssendung
populärwissenschaftliche Themen wie die Stringtheorie, Dimensionen der
Zeit, Illusion und Wirklichkeit, das Mögliche und Entscheidungen.

Der Schmetterlingseffekt wird durch ein gekochtes Ei ausgelöst, als ein
Regentropfen eine Telefonnummer verwischt, wodurch wichtige Verbindungen
verloren gehen.

Eine Version von Mr. Nobody erfährt aus einem aufgezeichneten Video von
einer älteren Version von ihm, dass er in den Erinnerungen eines Jungen
existiert, der sich nicht zwischen Vater und Mutter entscheiden kann.

Dem Journalisten, der das Sterben von Mr. Nobody dokumentiert, sagt er,
dass jedes seiner Leben das richtige Leben ist, jeder Pfad der richtige
Pfad, und dass diese Pfade vom sich entscheidenden Jungen erschaffen
werden.

Schließlich läuft die Zeit ab einem bestimmten Datum rückwärts auf den
Zeitpunkt der Entscheidung zwischen Vater und Mutter zu, doch der Junge
wählt nicht und läuft davon.

Interpretation: Mr. Nobody ist ein bemerkenswerter Film, der uns in die
Welt der Entscheidungen und ihrer Auswirkungen auf verschiedene
Realitäten entführt. Diese verschiedenen Realitäten wirken wie die Idee
vieler alternativer Welten, die der Form und Substanz nach einander
gleich sind und sich durch den Einfluss von alternativen Entscheidungen
der verschiedenen Versionen der handelnden Personen unterscheiden.

Es ist faszinierend zu sehen, wie der Protagonist, Mr. Nobody, seine
verschiedenen Lebenswege und Möglichkeiten durchlebt und wie dies unser
Verständnis von Entscheidungen und Realität herausfordert.

Hierbei entstehen parallel verschiedene Realitäten basierend auf
unterschiedlichen Entscheidungen und Verläufen. Jeder Weg ist dabei der
richtige Weg und real ist, was möglich ist. Deshalb ist die in der
Geschichte erzählte Geschichte einer Marsbesiedelung so real wie jede
andere Geschichte der Versionen von Mr. Nobody.

Die Konzeption der „Engel des Vergessens", die einem bei der Geburt
jegliches Wissen über mögliche Entwicklungen nehmen, kann als
metaphysische Metapher für diese Viele-Welten-Interpretation und den
Omegapunkt betrachtet werden. Und deshalb muss der Omegapunkt jede
Version von Mr. Nobody gelebt haben, damit er auch dieses Wissen hat. In
Mr. Nobody wird diese Idee in der Integration aller Geschichten und
Möglichkeiten dargestellt, die das Leben des Protagonisten ausmachen.

Persönliche Wertung: Unser Leben ist leidvoll (Geburt, Alter, Krankheit,
Sterben). Es rinnt uns wie Sand durch unsere Finger. Nicht einmal einen
Sinn, den wir ihm geben, können wir festhalten. Wir alle müssen sterben
und es wäre gut, dann ruhig und gelassen zu sein. Still ist unser Leben
selbst dann nicht, wenn alles ruhig ist und kein Geräusch zu hören ist.
Vollkommene Ruhe und Gelassenheit sind Zeitlosigkeit. Zeitlosigkeit ist
ein Zustand, den wir nie im Leben ganz erreichen können. Zeit ist
Zukunft, Gegenwart und Vergangenheit. Zukunft sind die offenen
Möglichkeiten und die zu ihnen gehörenden möglichen Geschichten.
Vergangenheit sind die Geschichten, an die wir uns erinnern können.
Gegenwart ist das bewusste Übergehen der Möglichkeiten und ihrer
möglichen Geschichten in die Vergangenheit. Zeitlosigkeit ist das
Verweilen in diesem Moment, der alle Möglichkeiten und ihre möglichen
Geschichten kennt und offen hält. An diesem Ort gibt es kein Leid, keine
Lust und kein Verlangen und das Wissen über alle Möglichkeiten und alle
möglichen Geschichten ist vollständig. An diesem Ort gibt es aber auch,
weil kein Leid und keine Lust, keine Zeit und kein Zeitverrinnen ist,
kein Ich, kein Selbst und kein Bewusstsein. Das ist nur dann kein
Widerspruch zu vollkommenem Wissen, wenn von diesem Ort aus
kontinuierlich alle Leben geboren, gelebt und erlitten werden, bis alles
kontinuierlich in die Zeitlosigkeit, in die Ruhe und Gelassenheit und in
vollkommenes Wissen eingeht.

\section{Das brandneue Testament}\label{das-brandneue-testament}

Das brandneue Testament von Jaco Van Dormael, 2015.

Autor: Jaco Van Dormael ist ein belgischer Regisseur und Drehbuchautor,
der wie Don DeLillo in einem katholisch geprägten Umfeld aufwuchs und
der für seine eigenwilligen, visuell wie gedanklich außergewöhnlichen
Filme bekannt ist. Geboren am 9. Februar 1957 in Ixelles, Belgien,
verbindet er in seinem Werk surreale Bildwelten mit philosophischen,
existenziellen und oft zutiefst menschlichen Fragestellungen. Seine
Filme entziehen sich einfachen Genrezuordnungen und kreisen um Zeit,
Möglichkeit, Freiheit, Leid und Sinn. Wie schon in Mr. Nobody arbeitet
Van Dormael auch in Das brandneue Testament mit Brüchen, Ironie und
metaphysischen Bildern, ohne dabei je den Blick für das Konkrete und
Alltägliche zu verlieren.

Handlung: Das brandneue Testament beginnt mit einer provokanten Setzung:
Gott existiert, und lebt in Brüssel. Nicht im Himmel, nicht im Jenseits,
sondern in einer Hochhauswohnung. Er ist kein liebender Schöpfer,
sondern ein cholerischer, kleinlicher und sadistischer Familienvater.
Seine Macht besteht nicht in Weisheit, sondern in Kontrolle. Er regiert
über Angst, Schuld und Strafe und verwaltet das Schicksal der Menschen
über einen Computer. Seine Tochter Éa lebt in dieser Wohnung wie in
einem Gefängnis. Der Vater tyrannisiert sie und ihre Mutter, der Sohn
Jesus ist nur noch als ferne, missachtete Erinnerung präsent. Als Éa den
verbotenen Raum betritt, in dem der Computer steht, wird sie brutal
bestraft. Diese Gewalt ist der Wendepunkt. Éa beschließt, sich dem Vater
zu entziehen. Sie sendet allen Menschen ihr jeweiliges Sterbedatum und
entzieht Gott damit sein wesentliches Machtinstrument: die Angst vor dem
Ungewissen. Anschließend flieht sie durch einen Waschsalon in die Welt
hinaus. Draußen sammelt sie sechs neue „Apostel", Menschen am Rand der
Gesellschaft, mit gebrochenen Biografien, unerfüllten Sehnsüchten und
einer oft wortlosen, aber tief verankerten Sinnorientierung. Parallel
dazu verliert Gott außerhalb seiner Wohnung rasch jede Autorität. In der
realen Welt ist er orientierungslos, aggressiv, unfähig zur Beziehung.
Er scheitert an Menschen, Institutionen und schließlich an sich selbst.
Am Ende wird er verhaftet und nach Usbekistan abgeschoben, nicht als
göttliche Strafe, sondern als groteske Konsequenz seiner Unfähigkeit,
Mensch unter Menschen zu sein. Währenddessen übernimmt die Mutter, lange
Zeit unsichtbar und unterdrückt, die Kontrolle über den Computer. Sie
widerruft die Sterbedaten und verwandelt die Welt. Gott wird entmachtet,
nicht durch Gewalt, sondern durch Bedeutungslosigkeit.

Interpretation: Auf den ersten Blick wird Das brandneue Testament häufig
als Religionssatire oder als Erzählung über die Gründung einer neuen
Religion gelesen. Diese Deutung greift jedoch zu kurz. Der Film handelt
nicht von Religion, sondern von einem bestimmten Gottesbild. Genauer:
von der Überwindung eines Gottes, der aus Angst, Schuld und Strafe
besteht. Dieser Gott trägt unverkennbar Züge eines
augustinisch-calvinistischen Denkens, ohne je explizit benannt zu
werden: totale Vorherbestimmung, moralische Kontrolle, Leiden als
notwendige Ordnung, Liebe als Belohnung für Gehorsam. Éas Handlung ist
keine Offenbarung und keine Heilslehre. Sie gründet nichts Neues. Sie
entzieht dem Vater lediglich die Macht. Indem die Menschen ihr
Sterbedatum kennen, verlieren sie nicht den Sinn, sondern die Angst. Sie
beginnen zu leben, nicht weil sie erlöst sind, sondern weil sie frei
werden. Die sechs Apostel verkörpern keine Dogmen, sondern gelebte
Möglichkeiten. Ihre Geschichten zeigen, dass Sinn nicht aus Ordnung
entsteht, sondern aus Beziehung, Aufmerksamkeit und Annahme. Die Mutter
schließlich ist die eigentliche Revolution des Films: Sie verändert die
Welt nicht durch neue Gebote, sondern durch das Aufheben der alten. In
diesem Sinn ist Das brandneue Testament keine Parodie des Christentums,
sondern eine leise Rückkehr zu dessen Geist, ohne ihn zu benennen. Nicht
Gesetz, sondern Befreiung. Nicht Schuld, sondern Beziehung. Nicht
Allmacht, sondern Ohnmacht, die akzeptiert wird.

Persönliche Wertung: Das eigentliche Problem des Films, und zugleich
seine größte Stärke, wird oft übersehen: Augustinus wird nicht erkannt.
Der Zuschauer lacht über Gott, über seine Lächerlichkeit, seine
Grausamkeit, seine Hilflosigkeit. Doch gemeint ist nicht Gott, sondern
ein Gottesbild, das tief in westlicher Kultur, Moral und
Selbstverständnis verankert ist. Der Film zeigt: Angst ist kein
Naturzustand. Schuld ist kein ontologisches Prinzip. Strafe ist keine
Voraussetzung für Sinn. Freiheit entsteht nicht durch neue Lehren,
sondern durch das Ende falscher Gewissheiten. Das brandneue Testament
ist deshalb kein Film über Religion, sondern über ihre Entmachtung dort,
wo sie zum Instrument der Kontrolle geworden ist. Er ist kein Aufruf zum
Glauben, sondern zur Menschlichkeit. Kein neues Testament, sondern das
Ende eines alten. Und vielleicht beginnt genau dort, wo Gott nicht mehr
herrscht, etwas, das man nicht benennen muss, um es zu leben.

\#\#\# Die andere Seite des Grenzwerts

Van Dormaels Filme sind nicht die Antwort auf DeLillo. Sie sind die
andere Seite desselben Grenzwerts. Wo DeLillo die Kohärenz zeigt, die
zerbricht, zeigt Van Dormael die Kohärenz, die sich vervielfacht. Wo
DeLillo die Stille als Grundzustand denkt, denkt Van Dormael sie als
Feld. Wo DeLillo den Omega-Punkt als Ausdünnung begreift, begreift Van
Dormael ihn als Fülle.

Es wäre ein Fehler, diese beiden Perspektiven gegeneinander
auszuspielen, oder sie zu einer Synthese zu bringen. Sie sind nicht
vereinbar. Sie sind gleichzeitig wahr, aber nicht zugleich sichtbar. Wer
von DeLillo kommt, sieht die Steine. Wer von Van Dormael kommt, sieht
die Verzweigungen. Beide sehen dasselbe, und sehen es anders. Der
Omega-Punkt ist nicht das eine oder das andere. Er ist der Rand, an dem
beide Perspektiven gelten, ohne sich zu berühren.

Das ist die eigentliche Entdeckung dieser Analyse: nicht die Ähnlichkeit
der Werke, sondern ihre Nicht-Vereinbarkeit. Nicht die gemeinsame
Botschaft, sondern der gemeinsame Grenzwert. Und vielleicht ist genau
das der Omega-Punkt, nicht ein Zustand, sondern ein Verhältnis. Nicht
ein Ziel, sondern ein Rand.

\section{Gott wohnt im dritten Stock}\label{gott-wohnt-im-dritten-stock}

„Gott wohnt im dritten Stock", Ein fiktiver Film Jaco Van Dormael
zugeschrieben.

Autor: Auch in seinem Film Gott wohnt im dritten Stock bleibt Jaco Van
Dormael dem Stil treu, der sein Werk seit jeher prägt: eine eigenwillige
Verbindung aus leiser Komik, existenzieller Schwere und philosophischer
Tiefenschärfe. Van Dormael erzählt nicht linear, nicht didaktisch,
sondern in Bildern, Brüchen und Bedeutungsverschiebungen. Wie schon in
Mr. Nobody und Das brandneue Testament interessiert ihn weniger die
Antwort als die Frage, weniger die Ordnung als das, was aus ihr
herausfällt. Dieser Film wirkt wie eine Verdichtung zentraler Motive
seines Œuvres: Zeit, Schuld, Freiheit, Möglichkeit, und die Frage, wie
Menschen leben, wenn die Instanz verschwindet, die vorgibt, alles zu
wissen.

Handlung: Gott wohnt im dritten Stock beginnt mit einer irritierenden
Klarstellung: Gott lebt in Brüssel. Nein, nicht Gott, sondern
Augustinus. Doch seine Tochter Éa nennt ihn so. Abfällig. Nicht aus
Trotz, sondern aus Erfahrung. Augustinus regiert seine Familie durch
Angst, Schuld und Strafe. Seine Wohnung ist kein Zuhause, sondern ein
System. Alles ist geregelt, alles bewertet. Der Sohn Michael hat dieses
System nicht überlebt. Sein Suizid ist eine Leerstelle, die nicht
betrauert, sondern verwaltet wird. Für Augustinus ist Michaels Tod ein
Scheitern der Ordnung. Für Éa ist er eine bleibende Präsenz. In Gedanken
spricht sie mit ihm, nicht als Jenseitsfigur, sondern als innerer Zeuge.
Éa verlässt die Wohnung mit einem Korb voller Wäsche. Der Waschsalon
wird zum Übergangsort. Sie kehrt nicht zurück. Draußen lebt sie unter
sechs Obdachlosen, die sich gegenseitig tragen, ohne Forderungen zu
stellen. Jeder von ihnen verkörpert eine andere Weise, Sinn zu leben:
durch Geduld, durch Rituale, durch Geschichten, durch Musik, durch
Aufmerksamkeit. Nichts davon wird erklärt, nichts wird missioniert.
Augustinus sucht seine Tochter. Doch außerhalb seiner Wohnung verliert
er jede Autorität. Seine Sprache greift ins Leere, seine Gewissheiten
werden nicht geteilt. Auch er gerät unter Obdachlose, aber er bleibt
fremd. Wo Éa angenommen wird, stößt er auf Distanz. Schließlich landet
er in einer Kirche, fällt einem Sozialarbeiter auf und wird in eine
psychiatrische Klinik eingewiesen. Parallel dazu beginnt die Mutter,
sich zu lösen. Sie sucht Éa, findet sie, ohne sie zurückzuholen.
Gemeinsam besuchen sie Augustinus in der Klinik. Das Gespräch bleibt
ruhig. Es gibt keine Versöhnung, keine Anklage. Nur Präsenz. Der Film
endet offen. Niemand wird erlöst. Niemand wird verurteilt. Doch Macht
hat sich verschoben.

Interpretation: Gott wohnt im dritten Stock ist kein Film über Religion.
Er ist ein Film über ein bestimmtes Verständnis von Gott, und über
dessen Scheitern. Augustinus steht für ein Weltbild, in dem Sinn durch
Ordnung entsteht, Schuld durch Abweichung, Sicherheit durch Kontrolle.
Dieses Weltbild ist nicht explizit religiös, aber tief religiös
strukturiert. Es braucht keine Dogmen, keine Rituale, keine heiligen
Texte. Es lebt von Angst. Éas Befreiung besteht nicht im
Glaubenswechsel, sondern im Verlassen dieses Systems. Die sechs
Obdachlosen repräsentieren keine neue Lehre, sondern gelebte
Alternativen: Sinn ohne Zwang, Gewissheit ohne Gewalt, Gemeinschaft ohne
Schuldrechnung. Ihre Haltungen erinnern an östliche Denkweisen, ohne
benannt zu werden. Sie sind nicht heilsbringend, sondern tragfähig.
Bemerkenswert ist, dass Augustinus nicht als Monster gezeichnet wird. Er
ist kein Bösewicht, sondern ein Mensch, der an seinen Gewissheiten
festhält, auch wenn sie niemandem mehr nützen. Seine Einweisung in die
Psychiatrie ist keine Strafe, sondern eine Unterbrechung. Ein Ort, an
dem sein Denken erstmals keine Wirkung mehr hat. Der Film zeigt: Gott
verschwindet nicht durch Widerlegung, sondern durch Irrelevanz. Macht
endet dort, wo Angst nicht mehr geteilt wird.

\section{\texorpdfstring{\hfill\break
}{ }}\label{section}

\section{Die dritte Perspektive: Entmachtung durch
Irrelevanz}\label{die-dritte-perspektive-entmachtung-durch-irrelevanz}

Es wäre verfehlt, Gott wohnt im dritten Stock einfach als eine weitere
Van-Dormael-Variation zu lesen. Der Film ist keine Wiederholung von Das
brandneue Testament. Er ist eine dritte Perspektive, eine, die weder
DeLillo noch Van Dormael ist.

Bei DeLillo verschwindet Gott durch Entzug. Die Stille ist leer. Bei Van
Dormael verschwindet Gott durch Bedeutungslosigkeit. Die Fülle ist da,
aber sie braucht Gott nicht mehr. Bei Gott wohnt im dritten Stock
verschwindet Gott durch Irrelevanz. Augustinus regiert, aber niemand
hört ihm zu. Seine Sprache greift ins Leere. Seine Gewissheiten werden
nicht geteilt. Er ist nicht widerlegt, er ist nicht mehr wichtig.

Das ist eine eigene Form der Entmachtung. Sie ist nicht metaphysisch.
Sie ist nicht theologisch. Sie ist sozial. Augustinus wird in die
Psychiatrie eingewiesen, weil sein Denken keine Wirkung mehr hat. Nicht
weil es falsch wäre, weil es keine Adressaten mehr findet. Das ist die
konsequenteste Form der Entmachtung: nicht Widerlegung, nicht
Bedeutungslosigkeit, sondern Verlust des Publikums.

Der Film zeigt damit etwas, das weder DeLillo noch Van Dormael zeigen:
dass das augustinische Denken nicht dadurch endet, dass man es
widerlegt. Es endet dadurch, dass niemand mehr hinhört. Und das ist kein
Triumph. Es ist eine stille, fast traurige Bewegung. Augustinus
verschwindet nicht in einem Feuerwerk. Er verschwindet in einem
Wartezimmer.

Persönliche Wertung: Dieser fiktive Film ist vielleicht Van Dormaels
konsequenteste Arbeit. Er ist leiser als Das brandneue Testament,
weniger verspielt als Mr. Nobody, aber existenziell schärfer. Das
eigentliche Drama liegt nicht in der Gewalt des Vaters, sondern in
seiner Unsichtbarkeit. Augustinus wird lange nicht erkannt, weder von
seiner Familie noch von sich selbst. Genau darin liegt die
gesellschaftliche Brisanz des Films: Er zeigt, wie tief Schuld- und
Angstlogiken verankert sind, selbst dort, wo niemand mehr von Gott
spricht. Gott wohnt im dritten Stock erzählt keine Erlösungsgeschichte.
Er erzählt von Entmachtung. Von der Befreiung aus einem Denken, das
Leben ordnet, aber nicht trägt. Am Ende bleibt Hoffnung, nicht, weil
alles gut wird, sondern weil nichts mehr gerechtfertigt werden muss. Und
vielleicht ist genau das der Moment, in dem Gott den dritten Stock
verlässt.

\section{Nachwort}\label{nachwort}

Dieser Essay hat nicht versucht, DeLillo und Van Dormael zu einer
gemeinsamen Lehre zu verbinden. Im Gegenteil: Je genauer man beide
liest, desto deutlicher wird ihre Nicht-Vereinbarkeit.

Bei DeLillo erscheint der Grenzbereich als Ausdünnung. Die Zeit verliert
ihre gewöhnliche Struktur, die Sprache verstummt, die Vermittlungen
verschwinden. Im Omega-Punkt liegt die Möglichkeit einer Zeitlosigkeit,
in der das Leben nicht mehr als fortlaufende Geschichte erscheint.

Bei Van Dormael, vor allem in \emph{Mr. Nobody}, erscheint derselbe
Grenzbereich als Fülle. Die Wirklichkeit verzweigt sich in
Möglichkeiten. Nicht eine Geschichte ist die wirkliche und alle anderen
sind falsch, sondern jede mögliche Geschichte besitzt ihre eigene
Wirklichkeit. Der Grenzbereich ist hier nicht Leere, sondern die
Gesamtheit der Möglichkeiten.

Beide Perspektiven lassen sich nicht zu einer dritten, harmonischen
Perspektive verschmelzen. Gerade deshalb ist der Omega-Punkt in diesem
Essay weder das eine noch das andere. Er ist der \textbf{Grenzwert}: der
unerreichbare Rand, an dem die Ausdünnung DeLillos und die Fülle Van
Dormaels zugleich denkbar werden, ohne einander aufzuheben.

In diesem Sinn liegt auch die entscheidende Differenz zu Teilhard de
Chardin. Teilhard macht aus dem Omega-Punkt einen Zielpunkt der
Entwicklung. Er gibt der kosmischen und geistigen Geschichte eine
Richtung, die auf Omega zuläuft. Der Grenzwert wird damit zum Telos.

\emph{Mr. Nobody} erlaubt eine andere Lesart. Wenn alle möglichen
Geschichten wirklich sind, dann ist Omega nicht das Ziel, das am Ende
einer Entwicklung erreicht wird. Es ist vielmehr die Grenze, an der die
Gesamtheit der Möglichkeiten gedacht werden muss. \textbf{Teilhard macht
aus dem Grenzwert eine Richtung; Van Dormael macht aus ihm einen
Möglichkeitsraum.}

Auch \emph{Das brandneue Testament} lässt sich von hier aus anders
lesen. Der Film entmachtet den strafenden, kontrollierenden Gott. Angst
verliert ihre religiöse Grundlage, und mit ihr soll die Freiheit des
Menschen zurückkehren. Doch damit ist das Problem noch nicht vollständig
gelöst. Denn ein Gottesbild kann seine äußere Macht verlieren und
dennoch als innere Struktur fortbestehen.

Genau hier setzt \emph{Gott wohnt im dritten Stock} an.

Augustinus ist in diesem Film nicht mehr der souveräne Herrscher. Seine
religiöse und soziale Macht ist gebrochen. Aber er ist deshalb nicht
geheilt. Was fortbesteht, ist die Struktur seines Denkens: Schuld,
Rechtfertigung, Kontrolle, Angst und das Bedürfnis nach einer letzten
Ordnung. Der Gott kann aus der Welt verschwinden und als innere Ordnung
im Menschen weiterleben.

Das ist die säkulare Form des Problems: \textbf{Der säkulare Mensch muss
nicht mehr an Gott glauben, um noch in Kategorien zu leben, die aus
einem bestimmten Gottesbild hervorgegangen sind.} Der äußere Herrscher
ist entmachtet; seine Krankheit ist geblieben.

Darum ist \emph{Gott wohnt im dritten Stock} keine bloße Fortsetzung von
\emph{Das brandneue Testament}. Während dort der Gott seine Macht
verliert, fragt der fiktive Film, was geschieht, wenn die Macht zwar
verschwunden ist, die von ihr geprägte Struktur aber weiterwirkt. Die
Entmachtung ist real. Die Überwindung ist es nicht notwendig.

Damit erhält auch die Spannung zwischen DeLillo und Van Dormael eine
weitere Dimension. Es geht nicht nur darum, was am Ende der Zeit
geschieht. Es geht auch darum, \textbf{welche Ordnungen wir auf dem Weg
dorthin mit uns tragen}.

Der Omega-Punkt ist deshalb keine Erlösungserzählung und keine neue
Metaphysik. Er ist auch kein Zustand, den der Mensch erreichen könnte.
Er ist der Rand, an dem die unterschiedlichen Möglichkeiten des Denkens
sichtbar werden:

die Stille und die Fülle,\\
die Ausdünnung und die Verzweigung,\\
die Zeitlosigkeit und die Möglichkeit,\\
die Entmachtung und das Fortleben einer alten Ordnung.

Nichts davon muss zu einer letzten Synthese werden.

Vielleicht liegt gerade darin die eigentliche Bedeutung des
Omega-Punktes: \textbf{Er ist nicht das Ende, das auf uns wartet,
sondern die Grenze unseres Denkens über Anfang, Ende, Zeit, Möglichkeit
und Sinn.}

Wer von DeLillo kommt, sieht an diesem Rand die Steine.

Wer von Van Dormael kommt, sieht die Verzweigungen.

Wer durch \emph{Das brandneue Testament} geht, sieht einen Gott, der
seine Macht verliert.

Wer \emph{Gott wohnt im dritten Stock} betritt, entdeckt, dass der
entmachtete Gott dennoch als innere Ordnung weiterleben kann.

Und das Pompeji-Projekt versucht nicht, eine dieser Perspektiven zur
letzten Wahrheit zu erklären. Es führt sie an denselben Rand.

Der Omega-Punkt ist deshalb \textbf{kein Ziel, keine Lösung und keine
Synthese}.

Er ist ein \textbf{unerreichbarer Grenzwert}.

Und vielleicht besteht seine Bedeutung gerade darin, dass er die
Spannung zwischen dem, was sich auflöst, und dem, was sich vervielfacht,
nicht beendet.

Er hält sie offen.

\end{document}
